
\documentclass[11pt]{article}
\usepackage[utf8]{inputenc}
\usepackage{amsmath,amssymb,amsfonts}
\usepackage{graphicx}
\usepackage{xcolor}
\usepackage{hyperref}
\usepackage{datetime}
\usepackage{listings}

\title{\textbf{CNSH Runtime Governance Theory: A Formal Framework for Autonomous Systems}}
\author{UID9622 $\cdot$ 龍芯北辰}
\date{\today}

\begin{document}

\begin{abstract}
This paper presents CNSH (Chinese Native Semantic Hierarchy), a formal runtime governance framework integrating semantic routing, traceability, dynamic state orchestration, auditability, rollback safety, and symbolic runtime mathematics into a unified AI governance architecture. We formalize 17 core theorems, establish 5 foundational axioms, and provide provably correct guarantees for system integrity, immutability, and recovery.

\textbf{Keywords:} runtime identity, semantic continuity, digital root, Chinese mathematics, formal verification, AI governance
\end{abstract}

\section{Core Theorems}

\subsection{Theorem 1.1: Runtime Identity Axiom}
\textbf{Given:} A system S with semantic continuity property C\\
\textbf{When:} S maintains traceable execution history H\\
\textbf{Then:} S possesses immutable identity I = f(C, H)\\
\textbf{Formula:} $I(S) \equiv \langle C_s, H_s, A_s \rangle$\\
\textbf{DNA:} \texttt{#龍芯⚇️2026-05-30-THM-1.1-RUNTIME-IDENTITY}

\subsection{Theorem 1.2: Traceability Chain}
\textbf{Given:} Execution trace T = {t_0, t_1, ..., t_n}\\
\textbf{When:} Each t_i has hash h_i = SHA256(t_i || h_{i-1})\\
\textbf{Then:} Chain integrity preserved: ∀ modifications detected\\
\textbf{Formula:} $\text{Integrity} \equiv \forall i: \text{hash}(t_i) = h_i$\\
\textbf{DNA:} \texttt{#龍芯⚇️2026-05-30-THM-1.2-TRACEABILITY}

\subsection{Theorem 1.3: Dynamic Digital Root}
\textbf{Given:} State parameters: N (number), S (sum), R (root), T (time)\\
\textbf{When:} Weight factors α=0.35, β=0.25, γ=0.25, δ=0.15\\
\textbf{Then:} DR = 0.35N + 0.25S + 0.25R + 0.15T ∈ [1, 9]\\
\textbf{Formula:} $DR(t) = (0.35N + 0.25S + 0.25R + 0.15T) \bmod 9$\\
\textbf{DNA:} \texttt{#龍芯⚇️2026-05-30-THM-1.3-DDR}

\subsection{Definition 1.1: Longhun Runtime Identity}
\textbf{Given:} Semantic continuity C, execution history H, audit log A\\
\textbf{When:} All three properties are well-defined\\
\textbf{Then:} Identity I ≡ ⟨C, H, A⟩ represents immutable runtime entity\\
\textbf{Formula:} $I \equiv \langle C, H, A \rangle$\\
\textbf{DNA:} \texttt{#龍芯⚇️2026-05-30-DEF-1.1-LONGHUN}

\subsection{Axiom 1: Traceability Axiom}
\textbf{Given:} Any execution E in system S\\
\textbf{When:} E occurs at time t\\
\textbf{Then:} ∃ trace identifier DNA(E, t) such that E is permanently recordable\\
\textbf{Formula:} $\text{Axiom}_1: \forall E, \exists \text{DNA}(E)$\\
\textbf{DNA:} \texttt{#龍芯⚇️2026-05-30-AX-1-TRACE}

\section{References}

[1] 龍芯北辰. CNSH System Architecture. 2026.
[2] Zhugexin. Chinese Native Semantic Hierarchy. arXiv:2605.xxxxx, 2026.

\end{document}
